% ============================================================
%  CellSim — Introduction and Related Work
%  Draft: 2026-04-29
%  Compile: pdflatex intro_related_work && bibtex intro_related_work
%            && pdflatex intro_related_work && pdflatex intro_related_work
% ============================================================
\documentclass[12pt,a4paper]{article}

\usepackage[utf8]{inputenc}
\usepackage[T1]{fontenc}
\usepackage[english]{babel}
\usepackage{amsmath,amssymb}
\usepackage[hidelinks]{hyperref}
% \usepackage{microtype}
\usepackage{geometry}
\usepackage{setspace}
\usepackage{parskip}

\geometry{left=3cm, right=3cm, top=3cm, bottom=3cm}
\onehalfspacing

\title{\textbf{CellSim: A Domain-Driven Software Framework\\
for Stochastic Tumour Evolution Simulation}\\[0.5em]
\large Introduction and Related Work (Draft)}
\author{}
\date{\today}

% ============================================================
\begin{document}
\maketitle
\tableofcontents
\newpage

% ============================================================
\section{Introduction}
\label{sec:introduction}
% ============================================================

Cancer is not a single disease but a heterogeneous family of conditions sharing a
common hallmark: the progressive and irreversible corruption of the cellular
mechanisms that govern proliferation, differentiation, and death
\cite{veneziano2025, chauviere2010}.
Tumour development is inherently a multi-scale, stochastic process that simultaneously
spans the molecular level---where individual DNA strand breaks, gene silencing events,
and protein–protein interactions occur---the cellular level---where
fitness differences between clones are resolved through competition---and the tissue
level---where the tumour microenvironment, immune surveillance, and spatial constraints
shape the evolutionary landscape \cite{su2024, hassan2024}.
This complexity is not merely descriptive: it has direct clinical consequences.
Interventions that succeed \emph{in vitro} routinely fail \emph{in vivo}, and
population-level epidemiological trends do not reliably predict individual-patient
outcomes, precisely because the stochastic nature of early oncogenesis is systematically
underestimated \cite{chauviere2010}.

In response to this challenge, the past six decades have witnessed an accelerating
effort to bring mathematical and computational tools to bear on cancer biology
\cite{pugh2025, chauviere2010}.
The aspiration is compelling: if cancer follows deterministic or stochastic laws at
the cellular level, then those laws should be amenable to formal analysis, and
models derived from them should generate testable, quantitative predictions.
The results, however, have been mixed.
As Chauviere and colleagues noted in their critical appraisal of the field,
\textit{``over-optimistic estimations about [mathematical oncology's] ability have
created unrealistic expectations''} \cite{chauviere2010}.
May articulated the underlying methodological trap with even greater precision:
models are frequently constructed with \textit{``an excruciating abundance of
detail in some aspects, whilst other important facets of the problem are misty''}
\cite{may2004}.
A model that is mathematically sophisticated yet biologically misspecified does not
merely fail to predict---it may actively mislead experimental design and clinical
decision-making.

The failure modes of incorrect models are not merely theoretical.
A deterministic ordinary differential equation (ODE) system may capture average
tumour growth dynamics while systematically missing the stochastic rare events---a
single cell acquiring a second loss-of-function allele in a tumour suppressor gene,
for instance---that are responsible for the transitions between disease phases
\cite{gillespie1977, moran1958}.
Cell–cell communication networks introduce frequency-dependent fitness effects
into the competitive landscape of the tissue: the selective advantage of an
immunosuppressive phenotype is not constant but depends on the current
composition of the tumour microenvironment \cite{su2024, nowak2004}.
These phenomena demand stochastic, agent-level modelling frameworks rather than
mean-field approximations.

Despite the recognised need, the translation from sound mathematical theory to
reproducible, extensible, and clinically applicable software remains incomplete.
The bioinformatics community has documented this gap extensively: technical debt,
lack of software engineering discipline, and absent continuous integration practices
compromise the reproducibility of computational findings and prevent their adoption
in clinical workflows \cite{ma2024}.
The problem is not a shortage of mathematical ideas---the theoretical foundations
are deep and well-established---but rather a shortage of software artefacts that
implement those ideas with the rigour, transparency, and modularity that
translational research demands.

This paper presents \textsc{CellSim}, a domain-driven, open-source software
framework for stochastic tumour evolution simulation, as a step towards closing
this gap.
\textsc{CellSim} is designed as a proof of concept for the intersection of rigorous
mathematical modelling, biologically grounded cellular mechanics, and
production-quality software engineering.
The remainder of this section motivates the three principal contributions of the
framework.
Section~\ref{sec:related_work} reviews the relevant literature across the three
axes---mathematical models, biological mechanisms, and software engineering
practice---and positions \textsc{CellSim} within that landscape.

% ============================================================
\section{Related Work}
\label{sec:related_work}
% ============================================================

\subsection{Mathematical Models in Cancer Biology}
\label{subsec:math_models}

Mathematical modelling in oncology has a history of over sixty years.
Bibliometric analysis confirms that the field has grown substantially in both volume
and interdisciplinary reach, with strong co-authorship networks spanning mathematics,
physics, and clinical oncology \cite{pugh2025}.
The diversity of modelling approaches reflects the multi-scale nature of the
underlying biology: deterministic models based on ordinary or partial differential
equations have been used to describe macroscopic tumour growth dynamics and
treatment response \cite{chauviere2010}; agent-based models (ABMs) capture
spatial heterogeneity and clonal competition at the cellular level; and
stochastic frameworks address the inherently probabilistic nature of molecular-level
events.

Among stochastic models, two foundational contributions stand out.
Moran's 1958 model of random processes in genetics \cite{moran1958} established
the mathematical framework for gene-frequency dynamics in a finite, fixed-size
population.
In Moran's formulation, instead of the simultaneous generational replacement assumed
by the Wright--Fisher model, births and deaths occur as individual, asynchronous
events: at each step, one individual is chosen uniformly at random to die and is
replaced by the offspring of a randomly sampled survivor.
This single modification---substituting individual stochastic replacement for
synchronous generation turnover---has profound consequences for evolutionary
dynamics: the rate of approach to homozygosity is exactly twice that of the
Wright model, because the random choice of the individual to die introduces
an additional source of variance \cite{moran1958}.
Translated to oncology, the Moran process provides a rigorous framework for
modelling clonal competition in homeostatic tissue, where cell number is
regulated and each division event is coupled to a death event.
The probability that a single mutant cell---say, one carrying a loss-of-function
allele in \textit{TP53}---ultimately fixes in the tissue, and the expected
time to fixation or loss, are directly computable from Moran's results.

Gillespie's Stochastic Simulation Algorithm (SSA) \cite{gillespie1977} addressed
a complementary problem: how to perform exact numerical simulation of coupled
stochastic reaction systems in continuous time.
Gillespie's key insight was philosophical as much as technical.
As he wrote, \textit{``the stochastic formulation of chemical kinetics has a
firmer physical basis than the deterministic formulation''}, because the deterministic
rate equation is merely the expected-value projection of the underlying master
equation \cite{gillespie1977}.
The SSA generates statistically exact trajectories of the master equation without
solving it analytically, making it possible to simulate rare molecular events---the
very events responsible for somatic mutation and phenotypic switching---with
mathematical exactness.
This is precisely the regime where ODE-based models are least reliable: the second
loss-of-function mutation in a tumour suppressor gene affecting a small number of
cells in a large tissue is a rare event that ODE approximations systematically
misrepresent.

When cellular fitness is not constant but depends on the current composition of
the population---as occurs in the tumour microenvironment, where immunosuppressive
cells alter the selective landscape for neighbouring cells \cite{su2024}---
the appropriate theoretical framework is evolutionary game theory.
Nowak and Sigmund demonstrated that evolutionary dynamics can be frequency-dependent
in a manner that is analytically tractable using replicator equations and evolutionary
stable strategy analysis \cite{nowak2004}.
The implication for tumour modelling is direct: the competitive advantage of a cell
with an immunoevasive phenotype is not fixed, but depends on the proportion of such
cells already present in the tissue, a feedback that purely kinetic models cannot
capture.
Nowak's comprehensive treatment of evolutionary dynamics \cite{nowak2006book}
further established that fixation probabilities and times in finite populations
under selection---central quantities for understanding clonal expansion---are
analytically accessible within these frameworks.

Despite this strong theoretical foundation, Hassan and colleagues \cite{hassan2024}
identify a persistent gap between mathematical models and clinical applicability:
most models are developed and validated on \emph{in silico} or cell-line data
without external validation against independent clinical cohorts.
The combination of machine learning and mechanistic mathematical modelling holds
promise for closing this gap, but the technical prerequisites---software
frameworks that are both mathematically rigorous and experimentally interoperable---
are still largely absent.

% ................................................................
\subsection{Bioinformatics, Software Engineering, and the
            Reproducibility Challenge}
\label{subsec:software}

A second body of literature relevant to this work concerns not the mathematical
content of models but their software realisation.
The rapid growth of biological data and the proliferation of computational tools have
created a secondary crisis: much of the software produced by researchers is
unreproducible, undocumented, and impossible to extend or validate independently
\cite{ma2024}.
Ma and colleagues \cite{ma2024} provide a systematic review of bioinformatics
software development practices, identifying technical debt---the accumulated cost
of shortcuts taken during development---as the primary obstacle to
reproducibility.
Their analysis identifies several critical deficiencies: absence of automated
testing, lack of version control discipline, hardcoded parameters that prevent
systematic sensitivity analysis, and monolithic architectures that resist
extension or combination with other tools.

Veneziano and colleagues \cite{veneziano2025} document the practical consequences
of this gap in the context of breast cancer bioinformatics: despite the
availability of rich transcriptomic, genomic, and proteomic data, the field
lacks software tools capable of integrating multi-omic signals into coherent
mechanistic narratives.
The gap is not conceptual but infrastructural.

Thomas and Barry \cite{thomas2003} argued over two decades ago for a
domain-oriented approach to software development in which software abstractions
are organised around the conceptual vocabulary of the domain---in this case,
cell biology and mathematical genetics---rather than around computational
convenience.
Their model-driven, domain-oriented programming (DOP) framework anticipates what
is now called Domain-Driven Design (DDD) \cite{thomas2003}: the principle that
the structure of the software should reflect the structure of the domain it models,
making the codebase intelligible to biologists, mathematicians, and software
engineers alike.
This approach is particularly relevant for scientific computing, where the primary
consumer of the software is often a domain expert who needs to reason about
model behaviour without reverse-engineering implementation details.

The intersection of the two deficits---missing mathematical rigour in
implementations, and missing software quality in mathematical tools---defines
the space that the present work aims to address.

% ................................................................
\subsection{Towards Integrated Frameworks: CellSim as a Proof of Concept}
\label{subsec:cellsim}

The preceding review identifies a three-way gap at the intersection of cancer
biology, stochastic mathematical modelling, and software engineering.
Mathematical oncology has produced powerful theoretical tools---the Moran process,
the Gillespie SSA, evolutionary game theory---but these tools remain largely
disconnected from software artefacts that implement them in a reproducible,
extensible, and biologically grounded way.
Pugh and colleagues \cite{pugh2025} confirm that interdisciplinarity is the
defining characteristic of successful mathematical oncology research, yet the
software infrastructure required to support multidisciplinary collaboration is
systematically underdeveloped.

\textsc{CellSim} is designed as a proof of concept for an integrated approach
to this problem.
It is an open-source, C++20 cellular evolution simulator that adopts a
hexagonal (ports-and-adapters) architecture \cite{thomas2003}, organising all
domain logic---cell lifecycle, genetic state, signal integration---in a pure
domain layer that depends on no external framework and is independently testable.

The biological model at the core of \textsc{CellSim} is grounded in established
cancer genetics.
Each \texttt{AgenticCell} carries structured representations of tumour suppressor
genes (\textit{TP53}, \textit{BRCA1}) with biologically constrained state
transitions: allelic states progress strictly from wild-type heterozygosity
(\texttt{+/-}) to homozygous loss (\texttt{-/-}), consistent with the two-hit
model of tumour suppressor inactivation \cite{veneziano2025}.
Two continuous damage accumulators---$D_1$ for DNA damage and $D_2$ for
immunosuppressive signalling---mediate the transition between six
biologically defined cell life stages (\texttt{BASELINE} $\to$ \texttt{UNSTABLE}
$\to$ \texttt{UNPROTECTED} $\to$ \texttt{PRIMER} $\to$ \texttt{TUMORAL}).
This staging mirrors the multi-step model of carcinogenesis documented in the
literature \cite{chauviere2010, veneziano2025}.

The cell lifecycle is implemented as a strictly ordered sequence of biological
decisions: intrinsic apoptosis mediated by \textit{TP53} is evaluated before
extrinsic immune-mediated apoptosis, which precedes neoplastic development and
cytoplasmic remodelling leading to division.
This ordering is not arbitrary but reflects the known hierarchy of cell fate
decisions in oncogenesis \cite{su2024}.

From an engineering perspective, \textsc{CellSim} implements continuous integration
with a dual-compiler matrix (GCC and Clang) and a comprehensive unit-test suite.
All architectural decisions are documented as Architecture Decision Records (ADRs),
providing a transparent audit trail that is essential for scientific reproducibility
and for onboarding contributors from different disciplinary backgrounds.

Critically, the architecture is designed from the outset for extension to the
mathematically rigorous stochastic frameworks reviewed in
Section~\ref{subsec:math_models}.
The Moran process extension will model tumour initiation in homeostatic tissue:
when a cell carrying a \textit{TP53} loss-of-function mutation divides and displaces
a wild-type neighbour, that is precisely a Moran replacement event, and the
probability of clonal fixation is directly given by Moran's formula \cite{moran1958}.
The Gillespie SSA extension will simulate rare somatic mutation events---including
the stochastic acquisition of the second \textit{BRCA1} hit---with mathematical
exactness, without the approximation errors that ODE-based models incur
in the low-copy-number regime \cite{gillespie1977}.
Tau-leaping \cite{gillespie1977} will provide a computationally tractable
approximation for the high-proliferation, large-cell-number regime associated
with the Big Bang dynamics of aggressive tumour growth.

Together, these components instantiate the three-way integration that the preceding
literature review identifies as necessary but currently absent: a biologically
grounded domain model, implemented with production-quality software engineering,
that can accommodate the full range of established stochastic mathematical
frameworks for tumour evolution.

% ============================================================
\newpage
\bibliographystyle{plain}
\bibliography{refs}

\end{document}
